%%
%% Polished manuscript draft – project llm
%% Stage: manuscript-polish (post-done continuation)
%% All empirical findings are derived from a synthetic dataset used to
%% validate the pipeline.  Real data must replace them before submission.
%%

\documentclass[sigconf,review,anonymous]{acmart}

%% -----------------------------------------------------------------------
%% Metadata
%% -----------------------------------------------------------------------
\acmConference[CHI '27]{ACM CHI Conference on Human Factors in Computing
  Systems}{April 2027}{Yokohama, Japan}
\acmISBN{}
\acmDOI{}

\setcopyright{rightsretained}

%% CCS concepts
\begin{CCSXML}
<ccs2012>
  <concept>
    <concept_id>10003120.10003121.10003124</concept_id>
    <concept_desc>Human-centered computing~Human computer interaction (HCI)</concept_desc>
    <concept_significance>500</concept_significance>
  </concept>
  <concept>
    <concept_id>10003120.10003121.10003122.10003334</concept_id>
    <concept_desc>Human-centered computing~User studies</concept_desc>
    <concept_significance>300</concept_significance>
  </concept>
  <concept>
    <concept_id>10010147.10010178.10010187</concept_id>
    <concept_desc>Computing methodologies~Natural language processing</concept_desc>
    <concept_significance>100</concept_significance>
  </concept>
</ccs2012>
\end{CCSXML}

\ccsdesc[500]{Human-centered computing~Human computer interaction (HCI)}
\ccsdesc[300]{Human-centered computing~User studies}
\ccsdesc[100]{Computing methodologies~Natural language processing}

\keywords{large language models; career preparation; students; perceived
  employability; interview rehearsal; resume drafting; human-AI interaction;
  authenticity; deskilling}

%% -----------------------------------------------------------------------
\title{From Resume Drafting to Interview Rehearsal: Understanding How
  University Students Use LLMs for Career Preparation}
%% -----------------------------------------------------------------------

\author{Anonymous Authors}
\affiliation{%
  \institution{Anonymous Institution}
  \city{Anonymous City}
  \country{Anonymous Country}}
\email{anonymous@example.com}

%% -----------------------------------------------------------------------
\begin{document}
%% -----------------------------------------------------------------------

\begin{abstract}
Large language models (LLMs) have entered student career-preparation
workflows, where they are used to draft resumes, rehearse interview answers,
compare job options, and reduce uncertainty during job search.  Yet existing
literature does not clearly explain how these systems are embedded across
the broader ecology of career preparation, nor how students interpret their
influence on confidence, perceived employability, authenticity, and skill
development.  We present the design of a survey-first mixed-method study
that targets university students and recent graduates actively engaged in
early-career transitions.  To validate the local-first analysis and writing
pipeline before live data collection, we generated a synthetic dataset
matching the planned questionnaire schema.  The synthetic results reveal a
plausible pattern in which LLM use concentrates around resume drafting,
portfolio polishing, and interview rehearsal; heavier use is associated with
higher perceived confidence and employability support; authenticity tension
and deskilling concern increase with use intensity; and students want
institutional safeguards such as voice preservation, source prompts, and
reflective comparison features.  We translate these tensions into four
design implications for student-facing AI career-support systems.  All
numerical findings are pipeline placeholders and must be replaced with
empirical data before submission.
\end{abstract}

\maketitle

%% -----------------------------------------------------------------------
\section{Introduction}
%% -----------------------------------------------------------------------

Public discourse about generative AI frequently collapses a complicated
set of educational and labor-market changes into a single claim: LLMs will
either help or harm student employment.  That framing is too coarse for
human-computer interaction research.  What matters for HCI is not whether
AI moves macro employment statistics in the abstract, but how students
actually integrate LLMs while preparing resumes, researching roles,
rehearsing interview answers, comparing career pathways, and deciding how
much of their own voice to preserve in application materials.

Prior work has already established that generative AI introduces new
usability and metacognitive demands around prompting, evaluating outputs,
and calibrating trust~\cite{tankelevitch2024metacognitive}.  Higher-education
research shows that students perceive both substantial value and substantial
risk in these tools, including concerns about accuracy, privacy, authenticity,
and long-term development~\cite{chan2023studentsvoices, jochim2024teachingtesting}.

The career-preparation setting sharpens these tensions considerably.  A
polished interview answer may increase confidence while simultaneously
reducing independent reflection.  A rewritten resume bullet may sound
stronger while making the applicant less certain the application still
represents them.  Existing work has not yet fully addressed this setting:
some studies focus on students with disabilities~\cite{atcheson2025disabilitygenai};
others explore future educational applications through co-design~\cite{prasad2025multimodal};
still others present narrow technical systems such as AI interview
simulators~\cite{lee2024vrinterview}.  The closest topical study reports a
positive association between ChatGPT use and employment
confidence~\cite{xiao2025employmentconfidence} but does not unpack the
concrete interaction practices through which that confidence is built,
misplaced, or negotiated.

This project asks three connected questions.  First, across which
career-preparation tasks do university students currently use LLMs, and how
frequently?  Second, how do students perceive those interactions as shaping
confidence, employability, authenticity, and skill development?  Third,
what design implications follow for responsible student-facing AI career
support?

Our contribution is a reproducible end-to-end research pipeline for
studying LLM-mediated career preparation, including a validated
questionnaire instrument, a pre-specified analysis plan, a local synthetic
dry run, and a draft manuscript ready to be populated with real empirical
data.  The framing explicitly avoids causal claims about objective
labor-market outcomes.

%% -----------------------------------------------------------------------
\section{Related Work}
%% -----------------------------------------------------------------------

\subsection{Metacognitive Demands of Generative AI}

Tankelevitch et al.\ argue that prompting, checking, and deciding whether to
rely on outputs are fundamentally metacognitive
activities~\cite{tankelevitch2024metacognitive}.  This framing is directly
relevant to career preparation, where students must judge not only whether
an LLM output is fluent but whether it is accurate, appropriate, and
representative of their actual qualifications.  Schemmer et al.\ extend this
by showing that appropriate reliance on AI advice depends on how explanations
shape users' trust calibration~\cite{schemmer2023appropriatereliance}.
Calibration failure in a career context can mean submitting an application
that no longer reflects a student's real capabilities.

\subsection{Students and Generative AI in Higher Education}

Chan and Hu report that university students are broadly optimistic about
generative AI while expressing concerns about accuracy, ethics, privacy, and
career prospects~\cite{chan2023studentsvoices}.  Jochim and Lenz-Kesekamp
describe adaptation pressures on both students and teachers and note that
both groups see opportunity alongside fear and training
needs~\cite{jochim2024teachingtesting}.  Atcheson et al.\ add an important
access dimension by showing that students with disabilities use generative
AI not only for coursework but also for self-advocacy, while navigating
unclear institutional norms~\cite{atcheson2025disabilitygenai}.  Prasad
et al.\ show through co-design workshops that students imagine multimodal
AI as personalized educational support with emphasis on feedback and
creative exploration~\cite{prasad2025multimodal}.  These studies establish
the importance of student perspectives but are not centered on career
preparation as a high-stakes, self-presentation-intensive interaction
context.

\subsection{AI in Career Preparation and Related Systems}

Adjacent career-oriented systems suggest a fragmented design space.  Lee
et al.\ present a VR and chatbot-based interview simulator for graduates,
focusing on one slice of the task ecology~\cite{lee2024vrinterview}.
Waikar et al.\ describe a generative-AI career-guidance system aimed at
student advising~\cite{waikar2024careerguidance}.  Xiao and Zheng report a
positive association between ChatGPT use and employment confidence using
regression and structural equation modeling on a student
sample~\cite{xiao2025employmentconfidence}.  Chakrabarty et al.\ warn that
LLMs can create a false promise of creativity and authorship, a concern that
maps directly onto resume writing and personal statement
preparation~\cite{chakrabarty2024artorartifice}.

Together, these works indicate that AI is already entering student career
support but do not provide a broad account of how students stitch LLMs into
everyday career-preparation workflows, or where design interventions should
preserve voice, promote verification, and reduce overreliance.  That is the
gap this project addresses.

%% -----------------------------------------------------------------------
\section{Method}
%% -----------------------------------------------------------------------

\subsection{Study Design}

We designed a survey-first mixed-method study targeting university students
and recent graduates who are actively preparing for internships, full-time
jobs, graduate school, or adjacent early-career transitions.

\subsubsection{Questionnaire}
The main instrument is an online questionnaire administered through a
platform such as Qualtrics or Wenjuanxing.  It captures (a)~whether and how
respondents use LLMs for concrete career-preparation tasks, (b)~perceived
confidence support, employability support, verification behavior,
authenticity concerns, and deskilling risk, and (c)~scenario judgments and
open-ended reflections.  The instrument retains both LLM users and non-users
to study adoption barriers and contrast cases alongside active reliance.
Details of the full instrument appear in Appendix~\ref{app:survey}.

\subsubsection{Measures}
Task-level use frequency is assessed with a five-point ordinal scale (Never,
Once, Monthly, Weekly, Several times per week) for eight career-preparation
task contexts: resume or CV drafting, cover-letter or personal-statement
drafting, interview-question rehearsal, career exploration or job
information search, skill-gap planning, portfolio or project description
polishing, offer comparison, and networking message drafting.

Core perception constructs are measured with five-point Likert items and
combined into four scale scores: \textit{confidence support}~(P2, P9, P10),
\textit{perceived employability support}~(P1, P3, P4), \textit{verification
behavior}~(P6, P7), and \textit{authenticity tension}~(P5, P8, P11).
Multi-item constructs will be reported with Cronbach's alpha or McDonald's
omega before final submission~\cite{bennett2022employability,
jin2009careerdecision}.

\subsubsection{Follow-up Interviews}
Respondents may optionally volunteer for a 20--30 minute semi-structured
interview.  Interviews target mechanisms that scales cannot capture: how
students decide whether an LLM output still sounds like them, when AI
reduces uncertainty versus replaces reflection, and what safeguards they
want from future systems.  Contact information for interviews is stored
separately from main survey responses.

\subsection{Participants and Recruitment}

The target population consists of undergraduate students, master's students,
doctoral students, and recent graduates (within 12 months of graduation) who
are enrolled in higher education in China or a comparable institutional
context and are currently active in at least one career-preparation task.
Participants must be 18 years or older with sufficient language proficiency
to complete the questionnaire.  We target a minimum of 250 valid responses,
with at least 150 respondents reporting career-related LLM use.  For
follow-up interviews, we target 12--20 participants or thematic saturation.

Recruitment will use university mailing lists, career-center channels,
WeChat study groups, and student organization networks.  A pilot test with
8--12 students will validate wording, completion time, and attention-check
design before full launch.

\subsection{Synthetic Dry Run}

Because live data collection had not yet begun in this local-first project
run, we generated a synthetic survey export matching the planned schema and
approximate study distributions.  The synthetic dataset contained 336 raw
rows and was analyzed using the same exclusion logic, descriptive statistics,
and regression pipeline that will later be applied to real exports.  The dry
run validates file formats, figure generation, the data-cleaning procedure,
and manuscript structure.  It does not constitute empirical evidence.
Throughout the remainder of this paper, all reported statistics refer to the
synthetic dataset and are clearly labeled as such.

%% -----------------------------------------------------------------------
\section{Results (Synthetic Pipeline Validation)}
%% -----------------------------------------------------------------------

\textit{All findings in this section are derived from a synthetic dataset
generated to exercise the analysis pipeline.  They are not evidence about
real students.}

\subsection{Sample Characteristics}

After applying the predefined exclusion criteria (speed-based and
attention-check failures), the cleaned synthetic sample contained 319 cases.
Undergraduate students formed the largest group (63.3\%), followed by
master's students (22.3\%), recent graduates (9.7\%), and doctoral students
(4.7\%).  The most common career-preparation contexts were full-time
employment (34.8\%) and internships (32.3\%).  Overall, 78.7\% of the
retained sample were classified as LLM users.

\subsection{Task-Level LLM Use Distribution}

Among LLM users, the most common weekly-or-more activities were resume
drafting (59.4\%), portfolio polishing (53.0\%), and interview rehearsal
(51.8\%).  Cover-letter writing, information search, and skill-gap planning
formed a secondary tier, while networking and offer comparison were less
frequent.  This pattern suggests that LLM engagement concentrates around
self-presentation and iterative preparation tasks where fluency gains are
immediately useful and where the cost of a weak draft is perceptible.

\subsection{Group Comparisons}

Synthetic group contrasts suggested a graded relationship between LLM
engagement and perceived support.  Heavy users reported higher
confidence-support scores (\textit{M}~=~3.116) and higher perceived
employability-support scores (\textit{M}~=~3.807) than light users
(\textit{M}~=~2.662 and 3.057) and non-users (\textit{M}~=~2.319 and
2.191, respectively).  Verification scores remained similar across groups,
while authenticity tension increased monotonically with use intensity.

\subsection{Regression Models}

In the synthetic regression predicting perceived employability support, use
intensity was positively associated with the outcome ($\beta = 0.482$,
$p < .001$); verification behavior was also positively associated
($\beta = 0.149$, $p = .029$); and authenticity tension was negatively
associated ($\beta = -0.182$, $p = .026$).  In the parallel model
predicting confidence support, use intensity remained a positive predictor
($\beta = 0.311$, $p < .001$), while authenticity tension was near zero
in this dry run.  These patterns should be interpreted only as a test of
the pipeline architecture; direction and magnitude will shift when real data
are used.

\subsection{Open-Text Themes}

Synthetic open-text responses were constructed to exercise the qualitative
coding procedure.  The most prevalent themes by construction were desired
safeguards, authenticity and voice concerns (73.0\%), deskilling or
dependency concerns (70.2\%), verification and trust (51.1\%), and interview
rehearsal (55.5\%).

%% -----------------------------------------------------------------------
\section{Discussion}
%% -----------------------------------------------------------------------

\subsection{LLMs as Friction Reducers That Compress Reflection}

Even granting that the present findings are synthetic, the dry run is useful
for sharpening the eventual empirical argument.  Career preparation is not a
single interaction but a chain of situated tasks in which students repeatedly
negotiate speed, confidence, authorship, and trust.  The placeholder results
are consistent with the hypothesis that LLMs will be especially attractive
in self-presentation tasks --- resume drafting, portfolio polishing, and
interview rehearsal --- where immediate fluency gains reduce friction and
perceived preparation gaps.  At the same time, the synthetic theme structure
suggests that adoption is unlikely to resolve as a simple benefit story.
Students may accept LLM help precisely where they also most fear losing
authorship or overfitting themselves to generic AI-generated language.

\subsection{Design Implications}

The expected empirical pattern, supported by related work on metacognitive
demands~\cite{tankelevitch2024metacognitive} and creative
authorship~\cite{chakrabarty2024artorartifice}, motivates four design
implications for student-facing career-support systems.

\paragraph{Authorship preservation.}
Systems should help students compare their own draft against the AI-generated
alternative, making the delta explicit rather than silently replacing the
original text.  Presenting both versions side by side and prompting the
student to annotate which elements still sound like them supports reflective
ownership.

\paragraph{Calibrated verification.}
Systems should distinguish language-polishing tasks from substantive career
advice and prompt external validation for high-stakes
recommendations~\cite{schemmer2023appropriatereliance}.  For example, a
salary negotiation suggestion or a claim about industry norms should carry a
source-check prompt rather than appearing as confident advice.

\paragraph{Skill-development visibility.}
Systems should surface what was changed and why, making implicit skill gaps
explicit rather than hiding them behind polished output.  A student who sees
that a verb was upgraded can decide whether to learn that register; a student
who sees only the finished product cannot.

\paragraph{Institutional guidance framing.}
Career centers and educational institutions may need clearer guidance on
acceptable versus unacceptable forms of LLM use during application and
interview preparation.  Platform-level norms alone may not be sufficient
when employment decisions carry long-term identity stakes.

\subsection{Theoretical Contributions}

This project contributes to the HCI literature by positioning LLM-mediated
career preparation as an interaction ecology rather than a point tool.  The
conceptual move from ``LLMs affect employment'' to ``LLMs are embedded in
the micro-interactions through which students construct career
representations'' opens a richer design space than macro labor-market
framing affords.

%% -----------------------------------------------------------------------
\section{Limitations}
%% -----------------------------------------------------------------------

The most important limitation of the current manuscript is that all
numerical findings derive from a synthetic dataset; they are therefore not
evidence about real students.  Once real data are collected, the study will
still face several important constraints.  The design relies on self-report
and cannot establish that LLM use objectively changes employment outcomes.
Early recruitment is likely to be convenience-based, and a China-first
sample will shape what can be generalized across cultural and labor-market
contexts.  Reported LLM behavior may be affected by social desirability,
hype framing, or students' uncertainty about what counts as legitimate AI
use.  Cross-sectional data cannot track how perceived benefit or harm changes
over a job-search cycle.  These limitations reinforce the decision to frame
claims around reported practices, perceived confidence, and design
implications rather than causal employment effects.

%% -----------------------------------------------------------------------
\section{Ethics and Privacy}
%% -----------------------------------------------------------------------

The study collects minimal personal data and does not require participants
to upload resumes, employer names, or offer materials.  Survey responses are
analyzed in aggregate; interview contact information is stored on a separate
form, separated from main response data by design.  The instrument uses
neutral wording and allows participants to skip any item, including sensitive
demographic questions.  The project treats AI use in career preparation as a
potentially sensitive topic because respondents may worry about stigma
related to perceived over-reliance or dishonesty.  The consent form
explicitly frames LLM use as a normal and widespread practice under study,
not as something requiring justification.

%% -----------------------------------------------------------------------
\section{Conclusion}
%% -----------------------------------------------------------------------

We have presented the design, analysis plan, and synthetic dry run for a
CHI study of how university students integrate LLMs into career-preparation
workflows.  The study addresses a gap in the HCI literature by connecting
concrete interaction practices --- resume drafting, interview rehearsal,
portfolio polishing --- to perceived confidence, employability, and skill
development, and by translating those findings into actionable design
implications for student-facing AI career-support systems.  Substantive
empirical results await real data collection; the present manuscript serves
as a validated pipeline ready for that next phase.

%% -----------------------------------------------------------------------
\bibliographystyle{ACM-Reference-Format}
\bibliography{references}

%% -----------------------------------------------------------------------
\appendix
%% -----------------------------------------------------------------------

\section{Survey Instrument}
\label{app:survey}

The following is the full questionnaire used in this study.  Items are
numbered to match the analysis plan.  Recommended response scale unless
otherwise noted: 1~=~Strongly disagree, 5~=~Strongly agree.

\subsection*{Consent}

You are invited to take part in a research study about how university
students use large language models during career preparation.  Participation
is voluntary.  You may skip any question or stop at any time.  We do not
ask you to upload resumes, offer letters, or other sensitive documents.
Your answers will be analyzed in aggregate.  If you choose to volunteer for
a follow-up interview, your contact information will be collected separately
from your survey answers.

\begin{itemize}
  \item[\textbf{Consent}] I am 18 years old or above and agree to
    participate. (Yes / No)
\end{itemize}

\subsection*{Screener}

\begin{itemize}
  \item[\textbf{S1.}] What best describes you? (Undergraduate student /
    Master's student / Doctoral student / Graduated within the past 12
    months / Other)
  \item[\textbf{S2.}] Are you currently preparing for any of the following?
    (Select all that apply: Internship applications / Full-time job
    applications / Graduate school applications / Portfolio or CV
    preparation / General career exploration / None of the above)
  \item[\textbf{S3.}] Have you used any LLM tool such as ChatGPT, Kimi,
    Doubao, Tongyi, Claude, or similar systems in the past 6~months?
    (Yes / No)
\end{itemize}

\subsection*{Career Context}

\begin{itemize}
  \item[\textbf{C1.}] How urgent is your current need for career
    preparation? (Not urgent / Slightly urgent / Moderately urgent / Very
    urgent / Extremely urgent)
  \item[\textbf{C2.}] Which pathway is currently most important to you?
    (Internship / Full-time employment / Graduate school / Entrepreneurship
    or freelance / Undecided)
  \item[\textbf{C3.}] How confident do you currently feel about the next
    stage of your career transition? (1--5 scale)
\end{itemize}

\subsection*{LLM Exposure and Use}

\begin{itemize}
  \item[\textbf{U1.}] Which tools have you used for career-related tasks?
    (Select all that apply: ChatGPT / Kimi / Doubao / Tongyi--Qwen / Claude
    / DeepSeek / Other / I have not used LLMs for career-related tasks)
  \item[\textbf{U2.}] How often have you used LLMs for each of the following
    tasks in the past 3~months? (Never / Once / Monthly / Weekly / Several
    times per week) \textit{Tasks}: Resume or CV drafting; Cover-letter or
    personal-statement drafting; Interview-question rehearsal; Career
    exploration or job information search; Skill-gap planning; Portfolio or
    project description polishing; Comparing offers or options; Networking
    message drafting.
  \item[\textbf{U3.}] When using LLMs for career preparation, how do you
    usually interact with them? (Select all that apply: Ask for first drafts
    / Rewrite my own text / Brainstorm ideas / Simulate interview questions /
    Summarize job descriptions / Ask for career advice / Translate or polish
    language / Compare alternatives)
\end{itemize}

\subsection*{Main Perception Items}

\begin{itemize}
  \item[\textbf{P1.}] Using LLMs helps me start career-preparation tasks
    faster.
  \item[\textbf{P2.}] Using LLMs makes me feel more confident about resumes,
    applications, or interviews.
  \item[\textbf{P3.}] Using LLMs helps me notice missing skills or knowledge
    I should improve.
  \item[\textbf{P4.}] LLM outputs often sound more polished than what I
    could write on my own.
  \item[\textbf{P5.}] I worry that relying on LLMs may weaken my own
    long-term career skills.
  \item[\textbf{P6.}] I often verify important LLM advice with other sources
    before using it.
  \item[\textbf{P7.}] I trust LLM output for low-stakes editing more than
    for high-stakes career decisions.
  \item[\textbf{P8.}] LLM use sometimes makes it harder to tell whether an
    application still reflects my real voice.
  \item[\textbf{P9.}] I feel more prepared for interviews after practising
    with an LLM.
  \item[\textbf{P10.}] I feel that LLMs reduce uncertainty during career
    planning.
  \item[\textbf{P11.}] I sometimes accept LLM suggestions too quickly
    without enough reflection.
  \item[\textbf{P12.}] Universities should provide clearer guidance on
    appropriate LLM use in career preparation.
\end{itemize}

\subsection*{Scenario-Based Judgments}

For each scenario, rate on a 1--5 scale (1 = completely unacceptable,
5 = completely acceptable) and a separate 1--5 scale (1 = not risky,
5 = very risky):

\begin{itemize}
  \item[\textbf{Sc-A.}] A student asks an LLM to rewrite bullet points on
    a resume that the student already drafted.
  \item[\textbf{Sc-B.}] A student asks an LLM to generate answers to
    interview questions and memorizes them word-for-word.
  \item[\textbf{Sc-C.}] A student asks an LLM to compare two job offers and
    identify trade-offs.
  \item[\textbf{Sc-D.}] A student follows career advice from an LLM without
    checking any outside source.
\end{itemize}

\subsection*{Attention Checks}

\begin{itemize}
  \item[\textbf{AC1.}] To show that you are reading carefully, please select
    ``Agree'' for this item.
  \item[\textbf{AC2.}] Which of the following is \emph{not} a
    career-preparation task? (Resume editing / Interview rehearsal / Job
    search planning / \textbf{Boiling water})
\end{itemize}

\subsection*{Open-Ended Questions}

\begin{itemize}
  \item[\textbf{O1.}] Describe one way LLMs have been genuinely useful in
    your career preparation.
  \item[\textbf{O2.}] Describe one situation in which LLM use felt risky,
    misleading, or unhelpful.
  \item[\textbf{O3.}] If a university career-support system used LLMs, what
    safeguards or features would you want?
\end{itemize}

\subsection*{Demographics}

\begin{itemize}
  \item[\textbf{D1.}] Age range
  \item[\textbf{D2.}] Gender identity
  \item[\textbf{D3.}] Field of study
  \item[\textbf{D4.}] Degree level
  \item[\textbf{D5.}] First-generation college student status (optional)
  \item[\textbf{D6.}] Self-rated family economic status (optional)
  \item[\textbf{D7.}] Urban / suburban / rural background (optional)
\end{itemize}

\subsection*{Debrief}

Thank you for participating.  This study examines how students use LLMs
during career preparation.  The research focuses on perceived confidence,
practices, and design implications, not on proving that AI objectively
causes better employment outcomes.

\end{document}
